%
% process has one thread with two stacks
%
\chapter{Scheduling}
\label{CH:SCHED}

Any operating system is likely to run with more processes than the
computer has CPUs, so a plan is needed to time-share the
CPUs among the processes. Ideally the sharing would be transparent to user
processes. A common approach is to provide each process
with the illusion that it has its own virtual CPU by
\indextext{multiplexing}
the processes onto the hardware CPUs.
This chapter explains how xv6++ achieves this multiplexing.

Before proceeding with this chapter, please read
(see \texttt{kernel/process.h}, \texttt{class process}),
(see \texttt{kernel/cpu\_manager.h}, \texttt{class cpu\_manager}),
(see \texttt{kernel/swtch.S}, \texttt{swtch}),
and the scheduling-related methods
(see \texttt{kernel/process\_manager.h}, \texttt{yield}),
(see \texttt{kernel/process\_manager.h}, \texttt{sched}),
and
(see \texttt{kernel/process\_manager.h}, \texttt{scheduler}).

%%
\section{Multiplexing}
%%

xv6++ multiplexes by switching each CPU from one process to
another in two situations. First, xv6++
switches when a process makes a system
call that blocks (has to wait), for example \lstinline{read} or
\lstinline{wait}.
Second, xv6++ periodically forces a switch to cope with
processes that compute for long periods without blocking.
The former are called voluntary switches,
the latter involuntary.

Implementing multiplexing poses a few challenges. First, how to switch from one
kernel thread to another?
The basic idea is to save and restore CPU registers,
though the fact that this cannot be expressed in C or C++ makes it tricky.
Second, how to force
switches in a way that is transparent to user processes?
xv6++ uses a hardware timer interrupt to drive context switches.
Third, all CPUs switch among the same set of processes, so a
locking plan is necessary to avoid mistakes such as two CPUs deciding
to run the same process at the same time.
Fourth, a process's
memory and other resources must be freed when the process exits,
but it cannot finish all of this itself.
Fifth, each CPU of a multi-core machine must remember which
process it is executing so that system calls affect the correct
process's kernel state.

%%
\section{Context switch overview}
%%

The term ``context switch'' refers to the steps involved in a CPU
leaving off execution of one kernel thread (usually for later
resumption), and resuming execution of a different kernel thread; this
switching is the heart of multiplexing.

xv6++ does not directly context
switch from one process's kernel thread to another process's kernel
thread; instead, a kernel thread gives up the CPU by context-switching
to that CPU's scheduler thread, and the scheduler thread picks a
different process's kernel thread to run, and context-switches to that
thread.

At a broader scope, the steps involved in switching from one user
process to another are:

1. A trap (system call or interrupt) from the old process's user space to its kernel thread  
2. A context switch to the current CPU's scheduler thread  
3. A context switch to a new process's kernel thread  
4. A trap return to the new user-level process  

%%
\section{Code: Context switching}
%%

The function \texttt{swtch()} in \texttt{kernel/swtch.S} contains the heart
of thread context switching: it saves the switched-from thread's CPU
registers, and restores the previously-saved registers of the
switched-to thread.

Each process object includes a saved context structure
(see \texttt{kernel/process.h}, \texttt{struct context})
that holds the thread's saved registers when it is not running.
A CPU's scheduler thread has its own saved context in the CPU structure
(see \texttt{kernel/cpu\_manager.h}, \texttt{struct cpu}).

When thread X wishes to switch to thread Y, thread X
calls \texttt{swtch(\&X\_context, \&Y\_context)}.
\texttt{swtch()}
saves the current CPU registers in X's context, then loads the contents
of Y's context into the CPU registers, then returns.

\begin{verbatim}
swtch:
        sd ra, 0(a0)
        sd sp, 8(a0)
        sd s0, 16(a0)
        ...
        sd s11, 104(a0)

        ld ra, 0(a1)
        ld sp, 8(a1)
        ld s0, 16(a1)
        ...
        ld s11, 104(a1)

        ret
\end{verbatim}

The \texttt{ret} returns to the instruction whose address is now in
\texttt{ra}, which was loaded from the restored context.
Thus control resumes exactly where the switched-to thread
previously called \texttt{swtch()}.

\texttt{swtch} saves callee-saved registers but not caller-saved registers.
The RISC-V calling convention guarantees that the caller preserves
caller-saved registers if needed.

%%
\section{Code: Scheduling}
%%

The scheduler exists as a special thread per CPU, each running
\texttt{scheduler}
(see \texttt{kernel/process\_manager.h}, \texttt{scheduler}).

A process that wants to give up the CPU must:

- Acquire its own process lock  
- Release any other locks  
- Change its state  
- Call \texttt{sched}  

You can see this sequence in
(see \texttt{kernel/process\_manager.cc}, \texttt{yield}),
(see \texttt{kernel/process\_manager.cc}, \texttt{sleep}),
and
(see \texttt{kernel/process\_manager.cc}, \texttt{kexit}).

\texttt{sched} calls \texttt{swtch} to save the current context and
switch to the scheduler context.

The scheduler loops over the process table
looking for a process whose state is \lstinline{RUNNABLE}.
When it finds one, it:

- Marks it \lstinline{RUNNING}
- Sets the per-CPU current-process pointer
- Calls \texttt{swtch} to run it

The process lock is held across the call to \texttt{swtch}.
This ensures that state changes and context saving occur atomically.

Releasing the lock too early could allow another CPU
to run a partially-saved process, causing stack corruption.

%%
\section{Code: mycpu and myproc}
%%

xv6++ must determine which CPU is currently executing.

While executing in the kernel, the RISC-V register
\lstinline{tp}
holds the CPU's hartid.

The function
(see \texttt{kernel/process\_manager.cc}, \texttt{mycpu})
uses \lstinline{tp}
to index the array of CPU structures.

Each CPU structure
(see \texttt{kernel/cpu\_manager.h}, \texttt{struct cpu})
contains:

- Pointer to current process  
- Scheduler context  
- Nested interrupt-disable count  

During boot,
(see \texttt{kernel/start.c}, \texttt{start})
initializes \lstinline{tp}.

When returning from user space,
(see \texttt{kernel/trampoline.S}, \texttt{uservec})
restores \lstinline{tp}.

The function
(see \texttt{kernel/process\_manager.cc}, \texttt{myproc})
returns the current process pointer safely by:

1. Disabling interrupts  
2. Calling \texttt{mycpu}  
3. Reading the process pointer  
4. Restoring interrupts  

%%
\section{Real world}
%%

The xv6++ scheduler implements a simple
\indextext{round robin}
policy.
Each runnable process gets a turn.

Real operating systems implement more sophisticated policies
involving priorities, fairness guarantees,
and throughput optimization.

%%
\section{Exercises}
%%

\begin{enumerate}

\item Modify xv6++ to switch directly from one process's kernel thread
to another without using the scheduler thread.
Be careful to prevent multiple CPUs from running the same process
and to avoid deadlocks.

\end{enumerate}
